% ROBOTICS COURSE PROJECT REPORT  --  Politecnico di Torino (PoliTo)
%  Hybrid KF + ANN Mobile-Robot Localization (after Yousuf & Kadri, 2020)
%
%  Compile:  pdflatex report.tex   (run it TWICE so the table of contents fills)
%
%  ADD AN IMAGE: replace a \figslot{<height>}{<caption>} line with, e.g.
%      \includegraphics[width=\linewidth]{figs/trajectory.png}
%  (keep the figure/caption around it). Put images in a ./figs/ folder.
%
%  EDIT NAMES/IDs: see the "EDIT ME" block just below \title.

\documentclass[11pt,a4paper]{article}

\usepackage[margin=2.5cm]{geometry}
\usepackage{amsmath,amssymb}
\usepackage{graphicx}
\usepackage{booktabs}
\usepackage{xcolor}
\usepackage{listings}
\usepackage{enumitem}
\usepackage{float}
\usepackage[hidelinks]{hyperref}

\definecolor{polired}{RGB}{140,20,30}
\definecolor{codebg}{RGB}{247,247,249}
\definecolor{codekw}{RGB}{0,90,160}
\definecolor{codecm}{RGB}{0,128,0}
\definecolor{codestr}{RGB}{170,55,55}

\lstdefinestyle{matlab}{
  language=Matlab, backgroundcolor=\color{codebg},
  basicstyle=\ttfamily\footnotesize,
  keywordstyle=\color{codekw}\bfseries,
  commentstyle=\color{codecm}\itshape,
  stringstyle=\color{codestr},
  numbers=left, numberstyle=\tiny\color{gray}, numbersep=7pt,
  breaklines=true, frame=single, rulecolor=\color{gray!40},
  showstringspaces=false, tabsize=2, columns=fullflexible,
  xleftmargin=14pt, framexleftmargin=14pt
}
\lstset{style=matlab}

% figure placeholder: \figslot{height}{caption}
\newcommand{\figslot}[2]{%
  \begin{figure}[H]\centering
  \setlength{\fboxsep}{0pt}%
  \fbox{\begin{minipage}[c][#1][c]{0.9\linewidth}\centering
     \itshape\color{gray}[ FIGURE PLACEHOLDER ]\\[2pt]
     \footnotesize replace with \texttt{\textbackslash includegraphics}
  \end{minipage}}\caption{#2}\end{figure}}

%==================== EDIT ME : title page details ============================
\title{\vspace{-1cm}\textbf{\color{polired}Sensor Fusion for Mobile-Robot Localization}\\[4pt]
\large Hybrid Kalman Filter and Artificial Neural Network Scheme\\
on the Pioneer P3-DX (CoppeliaSim)\\[8pt]
\normalsize Robotics Project --- Politecnico di Torino}

\author{Hussein Ghaddar \quad Fatima Zouhour \quad Rawan Elhalabi \quad Maya Kawssan\\[3pt]
\small Student IDs: \rule{2.5cm}{0.4pt} \quad \rule{2.5cm}{0.4pt} \quad \rule{2.5cm}{0.4pt} \quad \rule{2.5cm}{0.4pt}}
\date{\today}
%==============================================================================

\begin{document}
\maketitle
\thispagestyle{empty}

\begin{abstract}
\noindent We implement and study a multi-sensor localization scheme for a
differential-drive mobile robot, based on the method of Yousuf \& Kadri
(\textit{Robotica}, 2020). The robot is simulated in CoppeliaSim. Two Extended
Kalman Filters fuse GPS with inertial (KF-1) and odometric (KF-2) data; their
outputs are combined by a complementary filter. A feed-forward neural network is
trained, while GPS is available, to act as a pseudo-GPS during signal outages, and
its prediction is blended with KF-2 to recover position when GPS is lost. This
report describes the architecture, derives the filter equations we implemented,
presents results, and discusses the simplifications we made relative to the original
paper.
\end{abstract}

\tableofcontents
\vspace{4mm}

%==============================================================================
\section{Introduction}
%==============================================================================
Accurate self-localization is fundamental for autonomous mobile robots. No single
sensor is sufficient: GPS is accurate but unavailable indoors and in tunnels; inertial
(INS) and odometric dead reckoning are always available but drift over time through
integration error and, for odometry, wheel slippage. Fusing complementary sensors
gives a better estimate than any one alone.

This project reproduces, in a course setting, the localization scheme of Yousuf \&
Kadri (2020). The goal is to demonstrate understanding of (i)~Extended Kalman
filtering for nonlinear motion models, (ii)~complementary multi-sensor fusion, and
(iii)~the use of a neural network as a learned GPS substitute during outages. We do
not attempt a full reproduction of the paper; the scope and our simplifications are
stated explicitly in Section~\ref{sec:limitations}.

%==============================================================================
\section{System Architecture}
%==============================================================================
The pipeline is a chain of MATLAB scripts that pass data through \texttt{.mat} files:

\begin{enumerate}[nosep,leftmargin=2em]
  \item \texttt{BaseMatlab.m} --- drives the Pioneer P3-DX in CoppeliaSim and logs
        ground truth, GPS, IMU, odometry and a dead-reckoned INS.
  \item \texttt{ekf\_gps\_ins.m} --- 5-state EKF fusing GPS and INS (KF-1).
  \item \texttt{ekf\_gps\_odo.m} --- 3-state EKF fusing GPS and odometry (KF-2).
  \item \texttt{ekf\_fusion.m} --- complementary fusion of KF-1 and KF-2.
  \item \texttt{RobotANN\_Class.m} --- the feed-forward neural network.
  \item \texttt{ekf\_ANN\_fusion.m} --- trains the ANN and blends it with KF-2 during
        the GPS outage.
\end{enumerate}

The scripts must be run in this order, since each consumes the previous one's output.
A single runner script is given in Appendix~\ref{app:run}.

\figslot{6.5cm}{Block diagram of the fusion architecture: sensors feed KF-1 and KF-2;
their outputs are combined by the complementary filter while GPS is available; the
ANN, trained on this data, replaces GPS during an outage. (Redraw adapting the paper's
Fig.~1.)}

%==============================================================================
\section{Kalman Filter Models}
%==============================================================================
Each EKF runs the standard predict--correct cycle: the state is propagated through the
nonlinear motion model, then corrected by the GPS measurement \emph{only when GPS is
available}. During an outage the predicted state is carried forward (dead reckoning).
Both filters use the numerically stable \emph{Joseph form} for the covariance update
and re-symmetrise $P$ each step.

\subsection{KF-1: GPS + INS}
State $\mathbf{x}=[x,\,y,\,v_x,\,v_y,\,\theta]^\top$. Body-frame accelerations are
rotated into the world frame,
\[
a^w_x=\cos\theta\,a^b_x-\sin\theta\,a^b_y,\qquad
a^w_y=\sin\theta\,a^b_x+\cos\theta\,a^b_y,
\]
and the prediction is the constant-acceleration model
\[
\mathbf{x}^- =
\begin{bmatrix}
x+v_x\Delta t+\tfrac12 a^w_x\Delta t^2\\
y+v_y\Delta t+\tfrac12 a^w_y\Delta t^2\\
v_x+a^w_x\Delta t\\
v_y+a^w_y\Delta t\\
\theta+\dot\theta\,\Delta t
\end{bmatrix},\quad
F=\frac{\partial f}{\partial\mathbf{x}}=
\begin{bmatrix}
1&0&\Delta t&0&-\tfrac12\Delta t^2 a^w_y\\
0&1&0&\Delta t&\tfrac12\Delta t^2 a^w_x\\
0&0&1&0&-\Delta t\,a^w_y\\
0&0&0&1&\Delta t\,a^w_x\\
0&0&0&0&1
\end{bmatrix}.
\]
GPS observes position only, $H=[\,I_{2}\;\;0_{2\times3}\,]$.

\subsection{KF-2: GPS + odometry}
State $\mathbf{x}=[x,\,y,\,\theta]^\top$ with the unicycle model driven by forward
speed $v$ and yaw rate $\omega$ from the wheel encoders:
\[
\mathbf{x}^-=
\begin{bmatrix}
x+v\cos\theta\,\Delta t\\
y+v\sin\theta\,\Delta t\\
\theta+\omega\,\Delta t
\end{bmatrix},\quad
F=\begin{bmatrix}
1&0&-v\sin\theta\,\Delta t\\
0&1&\;\;v\cos\theta\,\Delta t\\
0&0&1
\end{bmatrix},\quad
G=\frac{\partial f}{\partial[v,\omega]}=
\begin{bmatrix}
\cos\theta\,\Delta t&0\\
\sin\theta\,\Delta t&0\\
0&\Delta t
\end{bmatrix}.
\]
Heading is wrapped to $(-\pi,\pi]$ with \verb|atan2(sin,cos)| before and after each
update. GPS again observes position, $H=[\,I_2\;\;0_{2\times1}\,]$.

\subsection{Complementary fusion}
The two position estimates are combined with a constant blending coefficient
$\alpha_1\in[0,1]$ (paper Eq.~2):
\[
[x_c,\,y_c]=\alpha_1\,[x_{kf1},\,y_{kf1}]+(1-\alpha_1)\,[x_{kf2},\,y_{kf2}].
\]
We used $\alpha_1=0.25$.

%==============================================================================
\section{Neural Network as Pseudo-GPS}
%==============================================================================
While GPS is available, a feed-forward network is trained to map a 10-dimensional
feature vector (acceleration and velocity magnitudes and their cumulative sums,
heading and cumulative heading, and odometric velocities) to the filtered position.
The architecture follows the paper: two hidden layers of 10 and 5 neurons with
log-sigmoid activations, a linear output layer, quasi-Newton (\texttt{trainoss})
training, input normalization (\texttt{mapminmax}), and weight-penalty regularization.

During a GPS outage the trained network produces a position prediction $[x_n,y_n]$,
which is blended with KF-2 (running on the last GPS update) by a coefficient
$\alpha_{1\_fuzzy}$ (paper Eq.~5):
\[
[x_{final},\,y_{final}]=\alpha_{1\_fuzzy}\,[x_n,\,y_n]
+(1-\alpha_{1\_fuzzy})\,[x_{kf2},\,y_{kf2}].
\]
In the original paper this coefficient is produced by a fuzzy logic system that
reacts to wheel slip; in our implementation we used a fixed value (see
Section~\ref{sec:limitations}).

%==============================================================================
\section{Results}
%==============================================================================
The figures and table below are placeholders to be filled from a single reproducible
run (fixed seed \texttt{rng(0)}, one GPS outage window). Save the plots into
\texttt{./figs/} and replace the placeholders.

\figslot{6.5cm}{Raw sensor data: ground truth vs.\ odometry vs.\ raw/smoothed GPS vs.\
INS (from \texttt{BaseMatlab.m}).}

\figslot{6.5cm}{KF-1 (GPS-INS) estimated trajectory against ground truth.}

\figslot{6.5cm}{KF-2 (GPS-odometry) estimated trajectory against ground truth.}

\figslot{6.5cm}{Final fusion during the GPS outage: truth, KF-2 alone, ANN alone, and
the ANN+KF-2 blend (from \texttt{ekf\_ANN\_fusion.m}).}

\begin{table}[H]\centering
\caption{Position RMSE (metres) --- to be filled from the run.}
\begin{tabular}{@{}lcc@{}}
\toprule
Estimator & RMSE during outage & RMSE whole run\\
\midrule
KF-1 (GPS-INS)        & \rule{2.2cm}{0.4pt} & \rule{2.2cm}{0.4pt}\\
KF-2 (GPS-odometry)   & \rule{2.2cm}{0.4pt} & \rule{2.2cm}{0.4pt}\\
Complementary fusion  & \rule{2.2cm}{0.4pt} & \rule{2.2cm}{0.4pt}\\
ANN alone             & \rule{2.2cm}{0.4pt} & \rule{2.2cm}{0.4pt}\\
ANN + KF-2            & \rule{2.2cm}{0.4pt} & \rule{2.2cm}{0.4pt}\\
\bottomrule
\end{tabular}
\end{table}

A short verification we recommend running is a heading-consistency check: plotting
KF-1's heading, KF-2's heading, and the course angle from ground truth on one axis
confirms that both filters and the simulator agree on the sign convention of the yaw
rate. If the two filter headings rotate in opposite directions, the complementary
fusion partially cancels and the sign in the INS/EKF heading update must be flipped.

%==============================================================================
\section{Simplifications and Limitations}\label{sec:limitations}
%==============================================================================
Relative to Yousuf \& Kadri (2020), we made the following simplifications, which we
note as honest scope decisions and possible future work:

\begin{itemize}[leftmargin=1.6em]
  \item \textbf{Reduced filter states.} We used a 5-state KF-1 and a 3-state KF-2,
        rather than the 11- and 7-state formulations of the paper. Our states capture
        the planar pose and velocity needed for 2-D localization.
  \item \textbf{Fixed blending coefficients.} Both fusion stages use constant weights.
        A natural improvement is inverse-variance weighting for the complementary
        filter (the EKFs already provide their covariances), which would automatically
        down-weight whichever filter is more uncertain during an outage.
  \item \textbf{Fuzzy logic system not implemented.} The paper's fuzzy slip-detector
        that sets $\alpha_{1\_fuzzy}$ was replaced by a fixed value. Implementing the
        three-rule Mamdani system (error between odometer and INS velocities, three
        triangular membership functions, centre-of-gravity defuzzification) is the
        main avenue for future work, and would allow the injected wheel slip to be
        compensated.
  \item \textbf{ANN training target and features.} We trained the network on the KF-2
        output and used the raw wrapped heading as a feature; training on the combined
        KF output and encoding heading as $(\sin\theta,\cos\theta)$ would likely
        improve generalization.
\end{itemize}

The core filtering and fusion pipeline is functional and the EKF formulations were
verified analytically. The above points are deliberate reductions in scope suited to
the course, not implementation defects.

%==============================================================================
\section{Conclusion}
%==============================================================================
We implemented a hybrid KF/ANN localization scheme for a simulated differential-drive
robot and reproduced its essential behaviour: two complementary EKFs fuse GPS with
inertial and odometric data, and a trained neural network substitutes for GPS during
outages. The project demonstrates the predict--correct EKF cycle on nonlinear motion
models, complementary multi-sensor fusion, and the pseudo-sensor role of a neural
network. Future work centres on adaptive (covariance- and fuzzy-based) blending and a
richer neural feature design.

\appendix
\section{Master run-script}\label{app:run}
\begin{lstlisting}[caption={\texttt{run\_all.m} --- runs the pipeline in the correct order.}]
clc; clear; close all;
run('BaseMatlab.m');     % -> kf_input.mat   (CoppeliaSim must be running)
run('ekf_gps_ins.m');    % -> kf1_output.mat
run('ekf_gps_odo.m');    % -> kf2_output.mat
run('ekf_fusion.m');     % complementary fusion
run('ekf_ANN_fusion.m'); % ANN + KF-2 during outage
disp('Pipeline complete.');
\end{lstlisting}

\section*{Reference}
S.~Yousuf and M.~B.~Kadri, ``Information Fusion of GPS, INS and Odometer Sensors for
Improving Localization Accuracy of Mobile Robots in Indoor and Outdoor
Applications,'' \textit{Robotica}, 2020. doi:10.1017/S0263574720000351.

\end{document}